\section{Stability Across Independent Training Runs}
\label{app:basis-stability}

This appendix documents the reproducibility evidence summarized in
\S\ref{subsec:stability} and the held-out grouping-recovery comparison
quoted in \S\ref{subsec:baseline-comparisons}. All experiments use
Qwen3.5-2B with three dictionaries per width configuration
(\(32\times\)/\(k{=}256\) and \(16\times\)/\(k{=}128\)), trained from
scratch with independent initialization, data-order, and sampling seeds;
held-out reconstruction is verified as equal across seeds (validation
top-1 agreement 0.805--0.848) before any stability comparison, so
differences below are not fidelity artifacts. All statements in this
appendix are \emph{within-recipe}: seeds vary, while width, sparsity,
and preprocessing are fixed. Cross-recipe variation is larger and is
disclosed at the end.

\subsection{Explanation-Level Stability}
\label{app:basis-stability-explanations}

The unit under test is the decomposition of a selected score: for each
of 63 curated contrasts (126 signed side-sets), we compare the token
sets that each seed's decomposition places on each side of the contrast.

\begin{table}[!tbp]
\caption{\textbf{Cross-seed stability of contrast explanations.}
Qwen3.5-2B, three seeds per width; means over seed pairs. Same-side
Jaccard compares side token-sets for the same contrast across seeds; the
cross-contrast null pairs decompositions of unrelated contrasts through
the identical pipeline.}
\label{tab:app-cross-seed-stability}
\centering
\footnotesize
\setlength{\tabcolsep}{5pt}
\renewcommand{\arraystretch}{1.08}
\begin{tabular}{@{}lcc@{}}
\toprule
Metric & \(32\times\)/\(k{=}256\) & \(16\times\)/\(k{=}128\) \\
\midrule
Same-side Jaccard (mean) & 0.214 & 0.240 \\
Cross-side leakage & 0.003 & 0.005 \\
Cross-contrast null & 0.014 & 0.018 \\
Contrasts above null p90 & 100\% & 100\% \\
Matched-projection correlation & 0.53--0.54 & 0.63--0.66 \\
Matched decoder cosine (used features) & 0.31--0.34 & 0.52--0.55 \\
\bottomrule
\end{tabular}
\end{table}

Individual atoms do not recur across seeds---matched decoder cosines
average \(\sim\)0.33 at \(32\times\)---consistent with the documented
seed-dependence of TopK dictionaries at scale
\citep{paulo2025stability}. Yet every tested contrast reproduces its
explanation above the null's 90th percentile, at
\(\sim\)15\(\times\) the null level and with negligible cross-side
leakage (\Cref{tab:app-cross-seed-stability}). Atom identity is not the
trustworthy unit of interpretation; the decomposition of a selected
score is.

\subsection{Feature-Group Matching}
\label{app:basis-stability-groups}

A complementary question is whether the feature-level token groups
themselves have an equivalent in another seed's dictionary, with indices
re-shuffled. The protocol was fixed before results: 100 features sampled
per seed pair from the bank-used set; each feature's group is its top-12
centered-row token set; the counterpart is the best-Jaccard match over
the full candidate dictionary; the null passes 500 frequency-matched
pseudo-groups through the identical best-of-\(D\) search, pricing in the
inflation of searching 65{,}536 candidates; recall@3 admits a greedy
union of at most three candidate features, testing whether instability
is feature splitting.

\begin{table}[!tbp]
\caption{\textbf{Cross-seed feature-group matching.} Qwen3.5-2B, ranges
over the three seed pairs per width. Null: frequency-matched
pseudo-groups through the identical search (best-single Jaccard median
0.059, p99 \(\approx\) 0.11; recall@3 0.25).}
\label{tab:app-feature-group-matching}
\centering
\footnotesize
\setlength{\tabcolsep}{4pt}
\renewcommand{\arraystretch}{1.08}
\begin{tabular}{@{}lcc@{}}
\toprule
Metric (range over pairs) & \(32\times\) & \(16\times\) \\
\midrule
Best-single Jaccard, median & 0.32--0.42 & 0.37--0.43 \\
Groups above null p99 & 89--90\% & 87--91\% \\
Strong equivalence (\(J\ge0.5\)) & 36--41\% & 40--47\% \\
Recall@3, median & 0.83--0.90 & 0.86--0.88 \\
Recall@3 above null p99 & 91--93\% & 92\% \\
Matched decoder cosine, med / p90 & 0.20--0.26 / 0.80--0.89 & 0.40--0.42 / 0.90--0.94 \\
\bottomrule
\end{tabular}
\end{table}

About 90\% of used feature groups have an above-chance counterpart in an
independently trained dictionary, roughly 40\% a near-exact single
match, and the median group is 83--90\% covered by a union of at most
three features (\Cref{tab:app-feature-group-matching}): token-group
structure reproduces, indexing does not, and the residual instability is
feature \emph{splitting} rather than noise. The decoder-cosine
distribution is bimodal---a stable core recurs almost exactly while the
rest recombine---and the narrower \(16\times\) atoms are more stable,
consistent with \citet{paulo2025stability}. For the below-null tail
(32/300 and 33/300 groups at the two widths), a direction-level check
finds decoder counterparts at cosine \(\ge 0.5\) for 5/32 and 10/33; a
qualitative pass suggests most of the remainder resemble the audit's
mixed/ambiguous class, but we have not rubric-scored that tail and do
not quote it as an audited rate. The strong-equivalence and recall@3
thresholds were fixed for this analysis rather than swept; we therefore
report them descriptively alongside the threshold-free medians.

\subsection{Held-Out Recovery of the Stable Grouping Core}
\label{app:basis-stability-loo}

The grouping comparison quoted in the main text asks which methods
recover reproducible grouping structure, on identical targets. For each
held-out seed, the \emph{core} is the set of grouped tokens stable in
both remaining seeds (leave-one-out construction; all methods evaluated
on identical cores; 95\% contrast-clustered bootstrap CIs).

\begin{table}[!tbp]
\caption{\textbf{Held-out recovery of the seed-stable grouping core.}
Qwen3.5-2B; recovery of the leave-one-out core by the held-out
same-recipe SRP dictionary, weighted row-kNN, and hard clustering, on
identical cores. The held-out SRP value (\(\approx\)0.75) is the
operative anchor: recall should be read against it, not against 1.0.}
\label{tab:app-loo-core-recovery}
\centering
\footnotesize
\setlength{\tabcolsep}{5pt}
\renewcommand{\arraystretch}{1.08}
\begin{tabular}{@{}lcc@{}}
\toprule
Method & \(32\times\) & \(16\times\) \\
\midrule
Held-out SRP (same recipe) & 0.749 [.73, .77] & 0.720 [.70, .74] \\
Weighted row-kNN & 0.493 [.44, .54] & 0.432 [.38, .48] \\
Hard clustering & 0.249 [.21, .29] & 0.206 [.17, .24] \\
\midrule
SRP/kNN ratio & 1.52 & 1.67 \\
Cross-contrast null & 0.014 & 0.018 \\
\bottomrule
\end{tabular}
\end{table}

Held-out SRP recovers 72--75\% of the stable core; row-kNN recovers
43--49\% and clustering 21--25\%, with SRP/kNN ratio CIs excluding
parity (\Cref{tab:app-loo-core-recovery}). Direct row geometry overlaps
the grouping structure well above the null but misses most of the
reproducible core that an independently trained same-recipe dictionary
recovers.

\subsection{Cross-Recipe Disclosure}
\label{app:basis-stability-cross-recipe}

The stability results above hold the training recipe fixed. Across
recipes the picture is weaker: the released paper dictionary---trained
under an earlier preprocessing recipe than the seed-variation
family---shares only 36\% of the seed-stable core and recovers 0.444 of
the leave-one-out cores, below row-kNN on that comparison. Cross-recipe
variation therefore exceeds within-recipe seed variation, and every
stability claim in this paper is scoped within-recipe. Comparisons in
this appendix use the seed-variation family end-to-end; no table mixes
dictionary families.
